% Stage 8 restored fragment — 2.1.3 Boundedness, compactness, rank deficiency, spectral compression (archive lines 1916-2204)
\subsection{Boundedness, compactness, rank deficiency, spectral compression}
\label{sec:ch213}

Spectral compression on compact support limits how sharply a finite aperture can beam \cite{jackson1999}. Sidelobe floors on a range trace the same law \cite{stratton1941}. Maxwell's retarded export to $\Gamma_R$ remains the governing constraint \cite{maxwell1873}.
The physical problem is observability under finite geometry: a compact source, a bounded
domain, and a finite aperture cannot couple to infinitely many independent field degrees of
freedom \cite{hansen1988}. Subsections~\ref{sec:ch211}--\ref{sec:ch212} fixed spaces and operator domains.
Subsection~\ref{sec:ch213} records boundedness and compactness properties, rank
deficiency of observation maps, and the viewpoint from which Section~\ref{sec:ch26}
develops spectral compression for radiation operators \cite{stratton1941}. The physical
content is familiar on a compact range: finite sources and finite probes cannot couple to
infinitely many independent far-zone degrees of freedom, so accessible rank is limited by
geometry and observation extent \cite{hansen1988}. Fredholm index computation
and complete spectral decomposition remain in Sections~\ref{sec:ch263}
and~\ref{sec:ch261}; here the constructive development establishes the lemmas on
boundedness, compactness, and rank deficiency those sections invoke by reference
\cite{stratton1941}.
In the display below, $\mathcal{F}_{\mathrm{rad}}\subset\mathcal{F}_\Omega$ denotes the
outgoing radiative subspace selected by $\Pi_{\mathrm{rad}}$; $\mathcal{S}_\omega$ and
$\mathcal{R}_\omega$ are the solution and radiation maps on $\mathcal{J}_\Omega$
\cite{harrington2001}.
The linear map representing \(T:\mathcal{X}\to\mathcal{Y}\) between Hilbert spaces is
\emph{bounded} if

\begin{equation}
\label{eq:ch2.operator-boundedness}
\|T\|_{\mathrm{op}}
=
\sup_{\|\mathbf{x}\|_{\mathcal{X}}=1}\|T\mathbf{x}\|_{\mathcal{Y}}
<\infty.
\end{equation}
Boundedness is equivalent to continuity at the origin and is the default
hypothesis for composition with measurement readouts \cite{stratton1941}. On the
Maxwell chain, operators split into bounded and unbounded classes. Trace and observation maps are bounded on declared spaces. The tangential trace
Eq.~\eqref{eq:ch2.trace-operator-bounded} satisfies
\(\|\gamma_\tau\|_{\mathrm{op}}<\infty\) from \(\HCurl(\Omega)\) to
\(H^{-1/2}(\Gamma)^{3}\) \cite{stratton1941}. Far-zone readouts from
Subsection~\ref{sec:ch201} define bounded functionals when restricted to
\(\mathcal{F}_{\mathrm{rad}}\):
\begin{equation}
\label{eq:ch2.far-observation-bounded}
\|\mathcal{O}_{\mathrm{far}}\|_{\mathrm{op}}
=
\sup_{\|\mathbf{u}\|_{\mathcal{U}}=1}
\big|\mathcal{O}_{\mathrm{far}}[\mathbf{u}]\big|
<\infty
\quad\text{on}\ \mathcal{F}_{\mathrm{rad}}.
\end{equation}
Laboratory maps \(\mathcal{O}_{\mathrm{meas}}\circ\mathcal{O}_{\mathrm{far}}\) in
Eq.~\eqref{eq:ch2.operational-measurement-chain} are bounded on the same subspace
because composition of bounded operators is bounded \cite{harrington2001}.
Volume and composite Maxwell operators are unbounded on full
\(\mathcal{U}_\omega\). The map \(\mathcal{M}_\omega\) of
Eq.~\eqref{eq:ch2.volume-operator-domain-range} has
\(\|\mathcal{M}_\omega\|_{\mathrm{op}}=\infty\) when
\(\mathcal{D}(\mathcal{M}_\omega)=\mathcal{U}_\omega\) because controlling
\(\|\mathcal{M}_\omega[\mathbf{u}]\|_{\mathrm{vol}}\) would require uniform curl
bounds beyond \(L^{2}\) control \cite{stratton1941}. Restriction to
\(\mathcal{D}(\mathcal{L}_\omega)\) with graph norm
Eq.~\eqref{eq:ch2.graph-norm} restores boundedness of the \emph{resolvent}
\(\mathcal{L}_\omega^{-1}\) on the physical branch where invertibility holds:
\begin{equation}
\label{eq:ch2.resolvent-bounded}
\|\mathcal{L}_\omega^{-1}\|_{\mathrm{op}}
<\infty
\quad\Longleftrightarrow\quad
\mathcal{L}_\omega\ \text{bijective on}\ \mathcal{D}(\mathcal{L}_\omega)/\ker\mathcal{L}_\omega.
\end{equation}
Equation~\eqref{eq:ch2.resolvent-bounded} specializes
Eq.~\eqref{eq:ch2.solution-map-stability}; exterior outgoing problems use
resolvent bounds on weighted subspaces rather than on all of
\(\mathcal{U}_\omega\) \cite{harrington2001}.
\begin{figure}[t]
\centering
\ChOneFigBlock{%
\begin{tikzpicture}[ch1fig, node distance=5mm and 7mm]
 \node[ch1box] (hc) {$\HCurl(\Omega)$};
 \node[ch1box, right=16mm of hc] (l2) {$L^{2}(\Omega)^{3}$};
 \node[ch1box, below=5mm of l2] (tr) {$H^{-1/2}(\Gamma)^{3}$};
 \draw[ch1flow] (hc.east) -- node[above,font=\scriptsize]{compact} (l2.west);
 \draw[ch1flow] (hc.south) -- node[right,font=\scriptsize]{$\gamma_\tau$} (tr.north);
\end{tikzpicture}%
}
\caption{Compact embedding \(\HCurl(\Omega)\hookrightarrow L^{2}(\Omega)^{3}\) on
bounded \(\Omega\); trace map is bounded but not compact
(Eq.~\eqref{eq:ch2.embedding-chain}).}
\label{fig:ch2.compact-embedding-schematic}
\end{figure}
~\eqref{eq:ch2.embedding-chain}. Rellich-type compactness for
\(\HCurl(\Omega)\hookrightarrow L^{2}(\Omega)^{3}\) on bounded domains implies
that bounded sequences in \(\HCurl\) have \(L^{2}\)-convergent subsequences
\cite{stratton1941}. Finite geometry cannot export arbitrarily fine angular complexity;
propagation therefore acts as a compression mechanism on electromagnetic structure.
An operator \(K:\mathcal{X}\to\mathcal{Y}\) is
\emph{compact} if bounded sets in \(\mathcal{X}\) map to relatively compact sets
in \(\mathcal{Y}\):
\begin{equation}
\label{eq:ch2.compact-operator-def}
K\ \text{compact}
\iff
K(B_{\mathcal{X}}(0,1))\ \text{has compact completion in}\ \mathcal{Y}.
\end{equation} 
Compact perturbations of identity produce Fredholm structure used in
Section~\ref{sec:ch263}; Subsection~\ref{sec:ch213} records the embedding lemmas
without computing index \cite{harrington2001}.
Radiation and Green maps on fixed bounded domains often decompose as
\begin{equation}
\label{eq:ch2.green-compact-split}
\mathcal{S}_\omega
=
G_{\mathrm{sing}}+K_{\mathrm{cpt}},
\end{equation}
with \(G_{\mathrm{sing}}\) capturing singular near-field kernels and
\(K_{\mathrm{cpt}}\) compact on \(\mathcal{J}_\Omega\) when sources have compact
support and \(\Omega\) is bounded \cite{stratton1941}. Equation
\eqref{eq:ch2.green-compact-split} is the operator template for
Section~\ref{sec:ch23}; kernel construction is not repeated here. The split
separates non-compact singular behavior at coincident points from compact
smoothing at finite separation \cite{harrington2001}. Rank deficiency enters through finite observation. An aperture-limited measurement
map \(\mathcal{O}_{\mathrm{meas}}:\mathcal{F}_{\mathrm{rad}}\to\mathbb{C}^{N}\)
has finite-dimensional range with \(N<\infty\):
\begin{equation}
\label{eq:ch2.finite-rank-observation}
\mathcal{O}_{\mathrm{meas}}
=
\sum_{n=1}^{N}\langle\cdot,\mathbf{f}_n\rangle_{\mathcal{U}}\,\mathbf{e}_n,
\end{equation}
for dual vectors \(\mathbf{f}_n\) and coordinate directions \(\mathbf{e}_n\).
Therefore \(\mathrm{rank}(\mathcal{O}_{\mathrm{meas}})\le N\) and
\begin{equation}
\label{eq:ch2.observation-null-space}
\ker\mathcal{O}_{\mathrm{meas}}
=
\left\{
\mathbf{u}\in\mathcal{F}_{\mathrm{rad}}:
\mathcal{O}_{\mathrm{meas}}[\mathbf{u}]=\mathbf{0}
\right\}
\end{equation}
is infinite-dimensional whenever \(\mathcal{F}_{\mathrm{rad}}\) is
infinite-dimensional \cite{stratton1941}. Equation
\eqref{eq:ch2.observation-null-space} is the operator form of the measurement
instrument rank limits in Section~\ref{sec:ch01.5}: finite modal sets cannot resolve all
angular content. Rank deficiency is structural, not a discretization artifact alone; the
spectral compression and singular-value decay law for radiation maps on compact support are
owned by Section~\ref{sec:ch26} \cite{stratton1941}.
\begin{figure}[t]
\centering
\ChOneFigBlock{%
\begin{tikzpicture}[ch1fig, node distance=5mm and 7mm]
 \node[ch1box] (f) {$\mathcal{F}_{\mathrm{rad}}$};
 \node[ch1box, right=7mm of f] (o) {$\mathcal{O}_{\mathrm{meas}}$};
 \node[ch1box, right=7mm of o] (cn) {$\mathbb{C}^{N}$};
 \node[ch1box, below=5mm of f] (ker) {$\ker\mathcal{O}_{\mathrm{meas}}$};
 \draw[ch1flow] (f) -- (o) -- (cn);
 \draw[ch1flow, dashed] (ker.south) -- (f.north);
\end{tikzpicture}%
}
\caption{Finite-rank infinite-dimensional field space maps to
\(\mathbb{C}^{N}\); \(\ker\mathcal{O}_{\mathrm{meas}}\) is infinite-dimensional for
\(N<\infty\).}
\label{fig:ch2.rank-deficient-observation}
\end{figure}
~\eqref{eq:ch2.radiation-operator-composition}.
Angular-spectrum reports that omit unresolved harmonics lie in the range of
\(\mathcal{O}_{\mathrm{meas}}\) but not in its injective sector
\cite{stratton1941}. Subsection~\ref{sec:ch213} does not classify report types;
Section~\ref{sec:ch01.5} remains authoritative. Singular values quantify rank deficiency. For compact
\(K:\mathcal{X}\to\mathcal{Y}\), the singular value sequence
\(\{\sigma_n(K)\}_{n\ge 1}\) decreases to zero, and
\begin{equation}
\label{eq:ch2.singular-value-decay}
\|K-R_N K\|_{\mathrm{op}}
=
\inf_{\mathrm{rank}\,R_N\le N}\|K-R_N\|_{\mathrm{op}}
\sim \sigma_{N+1}(K),
\end{equation}
where \(R_N\) ranges over rank-\(N\) operators \cite{stratton1941}. For radiation
maps truncated to \(N\) spherical or plane-wave modes, Eq.~\eqref{eq:ch2.singular-value-decay}
controls how much field content lies outside the retained modal subspace; explicit
VSH eigenbases are deferred to Section~\ref{sec:ch261} and Chapter~3 (\emph{Angular Momentum Operators and SO(3)})
\cite{harrington2001}. \emph{Spectral compression} names the reduction of infinite-dimensional radiation
dynamics to a finite or countable spectral model on the physical branch. Let
\(\{\phi_n\}\) be a spectral family in \(\mathcal{F}_\Omega\) and
\(\{P_n\}\) the associated one-dimensional projectors. A truncated radiation operator
\begin{equation}
\label{eq:ch2.spectral-compression-operator}
\mathcal{R}_\omega^{(N)}
=
\sum_{n=1}^{N}\lambda_n\,P_n\circ\mathcal{S}_\omega
\end{equation}
retains only modes with eigenvalues \(\lambda_n\) above a declared threshold
\cite{stratton1941}. Equation
\eqref{eq:ch2.spectral-compression-operator} is the abstract template for
Section~\ref{sec:ch26}; concrete eigenfunctions are not constructed here. On
bounded cavities, compact embeddings yield accumulating eigenvalues
\(\lambda_n\to\infty\) for interior resonances; on exterior domains, continuous
spectra require outgoing projectors rather than discrete sums
\cite{stratton1941}.
Compression error is bounded in operator norm when \(K\) is compact:
\begin{equation}
\label{eq:ch2.compression-error-bound}
\|\mathcal{R}_\omega-\mathcal{R}_\omega^{(N)}\|_{\mathrm{op}}
\le
C\,\sigma_{N+1}(K_{\mathrm{cpt}}),
\end{equation}
with \(K_{\mathrm{cpt}}\) from Eq.~\eqref{eq:ch2.green-compact-split} and
\(C\) independent of impressed data on bounded \(\mathcal{J}_\Omega\)
\cite{harrington2001}. Physically, Eq.~\eqref{eq:ch2.compression-error-bound}
states that truncating modal content beyond the \((N+1)\)st singular value cannot
be recovered by any rank-\(N\) surrogate operator; the quantitative decay law is
as in Section~\ref{sec:ch26} \cite{stratton1941}.
Rank deficiency of \(\mathcal{R}_\omega\) itself is distinct from observation rank
deficiency. Non-radiating sources satisfy
\(\mathcal{R}_\omega[\Jimp]=\mathbf{0}\) while
\(\mathcal{S}_\omega[\Jimp]\neq\mathbf{0}\) (Eq.~\eqref{eq:ch2.null-radiation});
Section~\ref{sec:ch272} develops that null space. Here the point is compositional:
\begin{equation}
\label{eq:ch2.rank-composition-law}
\mathrm{rank}(\mathcal{O}_{\mathrm{meas}}\mathcal{R}_\omega)
\le
\min\big\{\mathrm{rank}(\mathcal{O}_{\mathrm{meas}}),
\dim\mathcal{R}(\mathcal{R}_\omega)\big\},
\end{equation}
In Eq.~\eqref{eq:ch2.rank-composition-law},
\(\mathcal{O}_{\mathrm{meas}}\) is a finite-rank linear measurement functional
(pattern probe bank, modal fitter, or spherical-harmonic extractor) composed with
\(\mathcal{R}_\omega\); \(\mathrm{rank}(\cdot)\) is the dimension of the accessible
report subspace; and \(\mathcal{R}(\mathcal{R}_\omega)\) is the range of the radiation map
on \(\mathcal{J}_\Omega\).
So measurement rank caps reported angular content even when
\(\mathcal{R}_\omega\) is injective on admissible sources
\cite{stratton1941}. Both deficits must appear in uncertainty
budgets for angular-transport claims imported from Section~\ref{sec:ch01.2.3}. Fredholm structure note: on bounded \(\Omega\) with passive
\(\mathcal{B}_\Gamma\), \(\mathcal{L}_\omega\) is Fredholm on
\(\mathcal{D}(\mathcal{L}_\omega)\) modulo cavity kernels when written as
\begin{equation}
\label{eq:ch2.fredholm-split}
\mathcal{L}_\omega
=
\mathcal{L}_0+K,
\qquad
K\ \text{compact on graph domain},
\end{equation}
with \(\mathcal{L}_0\) an invertible reference operator on the quotient by
\(\ker\mathcal{L}_\omega\) \cite{stratton1941}. Index and solvability criteria are
developed in Section~\ref{sec:ch263}; Eq.~\eqref{eq:ch2.fredholm-split} supplies the
split Subsection~\ref{sec:ch212} stated the introductory relation in
Subsection~\ref{sec:ch202}.
As a worked example, a spherical near-field scanner records \(N=50\) tangential field samples on a
closed surface. The map \(\mathcal{O}_{\mathrm{meas}}\) has rank at most \(50\), so
angular harmonics beyond the resolved set lie in
\(\ker\mathcal{O}_{\mathrm{meas}}\). Applying Eq.~\eqref{eq:ch2.singular-value-decay}
to the compact smoothing operator from sample interpolation shows that recovered
far-zone coefficients error is controlled by \(\sigma_{51}\), not by Maxwell
residual alone \cite{stratton1941}. Increasing frequency without
increasing \(N\) lowers \(\sigma_{N+1}\) slowly for fixed aperture, worsening
compression error even when \(\mathcal{L}_\omega[\mathbf{u}_\omega]=\mathbf{0}\)
in the volume.
As a worked example, in a closed cavity, compact embedding yields discrete resonances and
Eq.~\eqref{eq:ch2.spectral-compression-operator} sums over cavity modes. In an
exterior problem with outgoing radiation condition, the same formula requires a continuous
spectral integral rather than a finite sum; using cavity compression on an
open antenna model conflates discrete storage modes with continuous radiation
modes \cite{stratton1941}. Section~\ref{sec:ch243} selects the outgoing branch;
Subsection~\ref{sec:ch213} only records that compression templates differ by domain class
\cite{stratton1941}.
Finite-geometry interpretation (compactness and rank). Compact embeddings and finite-rank observation maps express one enduring electromagnetic
truth: finite geometry destroys information. A bounded source support $\Omega_s$ cannot
excite arbitrarily fine angular structure at every observation point; a finite aperture
integrates over unresolved directions; a finite modal set truncates accessible rank.
Inaccessible angular structure is therefore not a numerical inconvenience but a consequence
of Maxwell propagation from compact sources through outgoing selection to finite
instruments. flux report is fundamentally limited, and the mathematics of compactness
makes that limit quantifiable through decaying singular values and infinite-dimensional
observation null spaces \cite{stratton1941}.
Subsection~\ref{sec:ch213} has established boundedness classes, compact embeddings, and
finite-rank observation structure. The quantitative singular-value decay for radiation maps
on compact support is developed in Section~\ref{sec:ch26}; here only the structural lemmas
are recorded. Subsection~\ref{sec:ch214}
completes the function-space foundation with self-adjoint realizations and admissibility
completions on bounded domains \cite{stratton1941}.
